% Options for packages loaded elsewhere
\PassOptionsToPackage{unicode}{hyperref}
\PassOptionsToPackage{hyphens}{url}
\documentclass[
  11pt,
]{article}
\usepackage{xcolor}
\usepackage[margin=1in]{geometry}
\usepackage{amsmath,amssymb}
\setcounter{secnumdepth}{-\maxdimen} % remove section numbering
\usepackage{iftex}
\ifPDFTeX
  \usepackage[T1]{fontenc}
  \usepackage[utf8]{inputenc}
  \usepackage{textcomp} % provide euro and other symbols
\else % if luatex or xetex
  \usepackage{unicode-math} % this also loads fontspec
  \defaultfontfeatures{Scale=MatchLowercase}
  \defaultfontfeatures[\rmfamily]{Ligatures=TeX,Scale=1}
\fi
\usepackage{lmodern}
\ifPDFTeX\else
  % xetex/luatex font selection
  \setmainfont[]{Cambria}
  \setmonofont[]{Consolas}
\fi
% Use upquote if available, for straight quotes in verbatim environments
\IfFileExists{upquote.sty}{\usepackage{upquote}}{}
\IfFileExists{microtype.sty}{% use microtype if available
  \usepackage[]{microtype}
  \UseMicrotypeSet[protrusion]{basicmath} % disable protrusion for tt fonts
}{}
\makeatletter
\@ifundefined{KOMAClassName}{% if non-KOMA class
  \IfFileExists{parskip.sty}{%
    \usepackage{parskip}
  }{% else
    \setlength{\parindent}{0pt}
    \setlength{\parskip}{6pt plus 2pt minus 1pt}}
}{% if KOMA class
  \KOMAoptions{parskip=half}}
\makeatother
\usepackage{longtable,booktabs,array}
\newcounter{none} % for unnumbered tables
\usepackage{calc} % for calculating minipage widths
% Correct order of tables after \paragraph or \subparagraph
\usepackage{etoolbox}
\makeatletter
\patchcmd\longtable{\par}{\if@noskipsec\mbox{}\fi\par}{}{}
\makeatother
% Allow footnotes in longtable head/foot
\IfFileExists{footnotehyper.sty}{\usepackage{footnotehyper}}{\usepackage{footnote}}
\makesavenoteenv{longtable}
\usepackage{graphicx}
\makeatletter
\newsavebox\pandoc@box
\newcommand*\pandocbounded[1]{% scales image to fit in text height/width
  \sbox\pandoc@box{#1}%
  \Gscale@div\@tempa{\textheight}{\dimexpr\ht\pandoc@box+\dp\pandoc@box\relax}%
  \Gscale@div\@tempb{\linewidth}{\wd\pandoc@box}%
  \ifdim\@tempb\p@<\@tempa\p@\let\@tempa\@tempb\fi% select the smaller of both
  \ifdim\@tempa\p@<\p@\scalebox{\@tempa}{\usebox\pandoc@box}%
  \else\usebox{\pandoc@box}%
  \fi%
}
% Set default figure placement to htbp
\def\fps@figure{htbp}
\makeatother
\setlength{\emergencystretch}{3em} % prevent overfull lines
\providecommand{\tightlist}{%
  \setlength{\itemsep}{0pt}\setlength{\parskip}{0pt}}
\usepackage{bookmark}
\IfFileExists{xurl.sty}{\usepackage{xurl}}{} % add URL line breaks if available
\urlstyle{same}
\hypersetup{
  hidelinks,
  pdfcreator={LaTeX via pandoc}}

\author{}
\date{}

\begin{document}

\section{The grokking exponent shrinks with task size: an empirical map
of α(K) with pre-registered robustness
checks}\label{the-grokking-exponent-shrinks-with-task-size-an-empirical-map-of-ux3b1k-with-pre-registered-robustness-checks}

\textbf{Pablo L. Rainieri Blasco}

Independent researcher --- Tucumán, Argentina

Draft v0.2 --- 2026-09-17 --- code and raw data:
\texttt{github.com/DoctorMilesBennettDyson/grokking-alpha-k}

\begin{quote}
\textbf{Draft status:} all data final, including the pre-registered
P1/P4 experiments (Section 5.3): P1 confirmed, P4 refuted. v0.2 =
submission candidate: related work completed with verified citations,
conclusion updated to the P1/P4 outcomes.
\end{quote}

\begin{center}\rule{0.5\linewidth}{0.5pt}\end{center}

\subsection{Abstract}\label{abstract}

Grokking --- delayed generalization long after training-set memorization
--- is known to depend strongly on weight decay (wd), with grokking time
following an approximate power law t\_grok ∝ wd\^{}α. We map how the
exponent α depends on task size on modular addition, training 2-layer
transformers across six moduli p ∈ \{31, 53, 67, 97, 113, 149\}
(training sets K from 480 to 11,100 examples) and five weight decays
over \textasciitilde200 runs on a single consumer GPU. We find that
\textbf{\textbar α\textbar{} shrinks monotonically-in-trend with K},
from -0.93 at K=480 to -0.58 at K=11,100: the larger the task, the less
weight decay accelerates grokking. We stress-test this map with
robustness checks pre-registered before execution: the effect is not an
initialization artifact (Kaiming/Xavier t\_grok ratio 1.28); grokking
time increases monotonically with batch size (1067 → 1314 → 1467 for bs
256 → 512 → 1024), indicating SGD noise \emph{accelerates} grokking; and
the power law survives a task-composition change (depth-2 sums, slope
-0.66) even though that regime exhibits no memorization phase at all ---
generalization there is ordinary, not grokked. Two negative results
bound the scope: plain SGD-momentum never groks at any tested setting
--- we identify the mechanism as weight-norm collapse under coupled L2 +
momentum (norm 42 → 0.02 within 1K steps) --- and an earlier theoretical
``pre-registration'' from our own program is disclosed as compromised
after a forensic reconstruction of its fit. The replacement predictions,
registered by public commit before their data collection began, are
reported at their 1-of-2 hit rate: the learning-rate direction is
confirmed, but batch-size invariance of α is refuted ---
\textbar α\textbar{} grows monotonically with batch size (0.64 → 0.73 →
0.80), so weight decay and batch size interact. All code, raw runs
(including failures), and pre-registration artifacts are public.

\subsection{1. Introduction}\label{introduction}

Grokking was reported by Power et al.~(2022): small transformers trained
on modular arithmetic memorize the training set quickly, then generalize
abruptly many thousands of steps later, with weight decay strongly
modulating the delay. Follow-up work explained the transition
mechanistically (Nanda et al., 2023: Fourier-feature circuits displacing
memorization) and as a phase transition (Liu et al., 2022). What remains
unmapped is the \textbf{interaction between regularization and task
size}: does weight decay accelerate grokking equally on small and large
tasks?

We answer empirically. Our contributions:

\begin{enumerate}
\def\labelenumi{\arabic{enumi}.}
\tightlist
\item
  \textbf{An α(K) map} (Section 3): the exponent of t\_grok ∝ wd\^{}α
  measured at six task sizes, N=3--8 seeds per cell, with all outliers
  and DNFs reported. \textbar α\textbar{} decreases from 0.93 to 0.58 as
  K grows 23×.
\item
  \textbf{Pre-registered robustness checks} (Section 4): initialization,
  batch size, optimizer, and task composition, with quantitative
  acceptance criteria fixed and dated before execution.
\item
  \textbf{Two mechanistic negative results} (Sections 4.2, 4.4):
  SGD-momentum weight-norm collapse, and the failure of a gradient-SNR
  reading of batch-size effects.
\item
  \textbf{A monotone weight-decay × batch-size interaction} (Section
  5.3): under a public, commit-timestamped pre-registration,
  \textbar α\textbar{} grows with batch size (0.64 → 0.73 → 0.80),
  refuting a separable-temperature model of the effect.
\item
  \textbf{A methods contribution} (Sections 5.2--5.3): we disclose, with
  the forensic reconstruction that exposed it, how a theory
  ``pre-registration'' inside our own program was contaminated by data
  it claimed to predate --- then re-run the confirmatory step correctly
  (commit-hash-timestamped, code-generated predictions), reporting its
  1-of-2 hit rate as the public commit obliges.
\end{enumerate}

Everything runs on one RTX 3070 (8 GB); total compute for all results is
under 150 GPU-hours.

\subsection{2. Setup}\label{setup}

\textbf{Task.} Modular addition \texttt{(a\ +\ b)\ mod\ p}, all p² pairs
enumerated, random 50/50 train/test split. Vocabulary: p digit tokens +
operator + equals. Depth-2 variant (Section 4.3):
\texttt{((a\ +\ b)\ +\ c)\ mod\ p}, subsampled to 500K examples.

\textbf{Model.} 2-layer transformer encoder, d\_model=128, 4 heads,
d\_ff=512, GELU, pre-norm; \textasciitilde423K parameters. Xavier init
unless stated.

\textbf{Training.} AdamW (β = 0.9/0.98), lr = 1e-3, 100-step warmup,
batch 512, 200K steps, full-batch evaluation every 100 steps. Weight
decay wd ∈ \{0.01, 0.1, 0.3, 1.0, 3.0\}.

\textbf{Definitions (fixed in pre-registration).} \texttt{t\_grok} =
first step with test accuracy ≥ 0.95 sustained for 3 consecutive
evaluations. A run is an \emph{outlier} if final test accuracy
\textless{} 0.95 or the sustained threshold trails the first crossing by
\textgreater{} 1000 steps; outliers are tabulated but excluded from cell
means. DNF = threshold never reached in budget; cells with DNFs are
lower bounds and excluded from slope fits.

\subsection{3. The α(K) map}\label{the-ux3b1k-map}

Master table (mean sustained t\_grok; full per-run tables
incl.~outliers/DNFs in \texttt{results/phase7/REPORT.md}):

{\def\LTcaptype{none} % do not increment counter
\begin{longtable}[]{@{}llllll@{}}
\toprule\noalign{}
p ~wd & 0.01 & 0.1 & 0.3 & 1.0 & 3.0 \\
\midrule\noalign{}
\endhead
\bottomrule\noalign{}
\endlastfoot
31 & 198,367 {[}3/8{]} & 39,660 {[}5/6{]} & 10,267 & 2,767 & 1,200 \\
53 & 153,700 {[}2/3{]} & 32,767 & 9,433 & 2,567 & 1,350 {[}2/3{]} \\
67 & 91,300 & 23,233 & 6,733 & 1,933 & 767 \\
97 & 30,133 & 9,267 & 4,100 & 1,314 {[}7/8{]} & 471 {[}7/8{]} \\
113 & 31,767 & 6,800 & 2,933 & 1,067 & 500 \\
149 & 13,233 & 5,767 & 2,333 & 833 & 617 {[}6/8{]} \\
\end{longtable}
}

The full grid is shown as a heatmap in \textbf{Figure 2}
(\texttt{fig2\_master\_heatmap.png}): t\_grok spans 471 to
\textasciitilde198,000 steps and falls smoothly along both axes.

Log--log slope of t\_grok vs wd per modulus (\textbf{Figure 1},
\texttt{fig1\_alpha\_vs\_K.png}):

{\def\LTcaptype{none} % do not increment counter
\begin{longtable}[]{@{}lll@{}}
\toprule\noalign{}
p & K & α \\
\midrule\noalign{}
\endhead
\bottomrule\noalign{}
\endlastfoot
31 & 480 & \textbf{-0.926} \\
53 & 1404 & -0.866 \\
67 & 2244 & -0.861 \\
97 & 4704 & -0.731 \\
113 & 6384 & -0.736 \\
149 & 11100 & \textbf{-0.577} \\
\end{longtable}
}

\begin{figure}
\centering
\pandocbounded{\includegraphics[keepaspectratio,alt={\textbar α\textbar{} --- the exponent of t\_grok ∝ wd\^{}α --- against training-set size K (log scale), one point per modulus. \textbar α\textbar{} falls from 0.93 (p=31) to 0.58 (p=149); the two plateaus (53→67, 97→113) are inside seed noise.}]{fig1_alpha_vs_K.png}}
\caption{\textbar α\textbar{} --- the exponent of t\_grok ∝ wd\^{}α ---
against training-set size K (log scale), one point per modulus.
\textbar α\textbar{} falls from 0.93 (p=31) to 0.58 (p=149); the two
plateaus (53→67, 97→113) are inside seed noise.}
\end{figure}

\begin{figure}
\centering
\pandocbounded{\includegraphics[keepaspectratio,alt={Master grid: mean sustained t\_grok over six moduli × five weight decays (\textasciitilde200 runs). t\_grok spans 471 to \textasciitilde198,000 steps and falls smoothly along both axes.}]{fig2_master_heatmap.png}}
\caption{Master grid: mean sustained t\_grok over six moduli × five
weight decays (\textasciitilde200 runs). t\_grok spans 471 to
\textasciitilde198,000 steps and falls smoothly along both axes.}
\end{figure}

Two honest observations. First, the \emph{trend} is clearly decreasing
(\textbar α\textbar{} falls 0.35 over a 23× range in K), but strict
monotonicity is not resolvable at N=3/cell: there are two local plateaus
(53→67 and 97→113, both \textbar Δ\textbar{} = 0.005, far inside noise).
Second, the noisiest cell in the grid (p=31, wd=0.01: 3/8 clean, the
rest DNF or unstable) anchors the steep end; its slope moved from -0.86
to -0.93 when five seeds were added. Both facts carry into every
downstream fit.

Interpretation, informally: weight decay is a \emph{circuit-complexity
price}. On small tasks (few constraints) many circuits fit the data and
wd is decisive in forcing the simple, generalizing one --- so t\_grok
responds strongly. On large tasks the data itself already excludes most
circuits, and wd has less selection left to do.

\subsection{4. Pre-registered robustness (Fase
II)}\label{pre-registered-robustness-fase-ii}

All four hypotheses, acceptance intervals, and analysis rules were fixed
and dated 2026-04-27, before execution began 2026-04-28
(\texttt{prereg/PRE\_REGISTRATION\_FASE\_II.md}). Execution completed
2026-07-07 after a crash-and-resume documented in that file's changelog.
Verdicts:

\subsubsection{4.1 H3 --- Initialization
(CONFIRMED)}\label{h3-initialization-confirmed}

Criterion: Kaiming/Xavier sustained-t\_grok ratio ∈ {[}0.5, 1.5{]} at
the p=97, wd=1.0 baseline, 5 fresh seeds each. Observed: \textbf{1.28}
(1640 ± 288 vs 1280 ± 84). The Xavier arm doubles as an independent-seed
replication of the baseline (1280 vs 1314 ± 195 from earlier phases).

\subsubsection{4.2 H2 --- Optimizer (negative result;
scope-bounding)}\label{h2-optimizer-negative-result-scope-bounding}

Criterion: SGD-momentum (lr 1e-2, momentum 0.9) preserves the sign of α.
Observed: \textbf{0/15 runs grok --- all end at chance} (train loss
pinned at the uniform-distribution value). A differential diagnostic
isolates the mechanism: with wd=0, the same SGD path memorizes the
training set in \textasciitilde1K steps (code is healthy); with wd=1.0
the weight norm collapses 42 → 0.02 within 1K steps and the network
emits uniform logits permanently. Decoupled SGDW at wd=1.0 collapses
similarly. AdamW survives because adaptive normalization keeps update
magnitudes O(lr) as weights shrink, while SGD's gradients shrink
\emph{with} the weights, so decay always wins. \textbf{The α(K) map is
therefore a statement about AdamW-class optimizers}; ``weight decay'' is
not a single mechanism across optimizers, echoing the decoupling
analysis of Loshchilov \& Hutter (2019).

\subsubsection{4.3 H1 --- Task composition (confirmed numerically;
regime boundary
found)}\label{h1-task-composition-confirmed-numerically-regime-boundary-found}

Criterion: at depth-2 (p=97), slope ∈ {[}-0.94, -0.54{]}. Observed:
\textbf{-0.658} over the three wd cells that generalize (wd ≥ 1.0: 0/3
reach criterion). But in every clean depth-2 run, train and test
accuracy rise \emph{together} ---
\texttt{t\_train\_saturated\ ==\ t\_grok} in all cells --- because with
250K training examples and 423K parameters, memorization is impossible.
\textbf{Figure 3} contrasts the two regimes directly: a canonical
grokking gap at depth-1 (train saturates \textasciitilde300 steps, test
generalizes \textasciitilde25,000 steps later) beside the depth-2 curve
where train and test are indistinguishable throughout. \textbf{There is
no grokking at depth-2}; the wd power law survives in a regime of
ordinary generalization. We report this as a boundary of the grokking
regime (the memorization phase requires K small relative to capacity),
and explicitly \emph{not} as ``grokking is composition-invariant''.

\begin{figure}
\centering
\pandocbounded{\includegraphics[keepaspectratio,alt={Left: canonical grokking at depth-1 (p=97, wd=0.01) --- train accuracy saturates by \textasciitilde300 steps, test accuracy generalizes \textasciitilde25,000 steps later. Right: depth-2 (p=97, wd=0.1) --- train and test rise together; there is no memorization phase and therefore no grokking, yet the wd power law still holds (§4.3).}]{fig3_regime_contrast.png}}
\caption{Left: canonical grokking at depth-1 (p=97, wd=0.01) --- train
accuracy saturates by \textasciitilde300 steps, test accuracy
generalizes \textasciitilde25,000 steps later. Right: depth-2 (p=97,
wd=0.1) --- train and test rise together; there is no memorization phase
and therefore no grokking, yet the wd power law still holds (§4.3).}
\end{figure}

\subsubsection{4.4 H4 --- Batch size (CONFIRMED; kills a naive theory
reading)}\label{h4-batch-size-confirmed-kills-a-naive-theory-reading}

Criterion: sustained t\_grok monotonically increasing with batch size at
p=97, wd=1.0. Observed: \textbf{1067 (bs=256) → 1314 (bs=512, reference)
→ 1467 (bs=1024)} --- monotone as pre-registered. The quantitative
sub-prediction (each doubling +20--80\%) held for the first doubling
(+23\%) and failed for the second (+12\%). The direction matters more
than the magnitudes: if minibatch noise limited \emph{signal} (a
gradient-SNR reading), larger batches should grok \emph{faster}. They
grok slower. Minibatch noise \textbf{accelerates} grokking, consistent
with noise-as-exploration accounts of the memorization-to-generalization
transition. \textbf{Figure 4} shows the batch-size effect (left) and the
SGD weight-collapse mechanism (right: AdamW train loss crashes to zero
while SGD stays pinned at the uniform-distribution value ln(99)).

\begin{figure}
\centering
\pandocbounded{\includegraphics[keepaspectratio,alt={Left: sustained t\_grok increases monotonically with batch size at p=97, wd=1.0 (1067 → 1314 → 1467 for bs 256 → 512 → 1024): minibatch noise accelerates grokking. Right: the SGD-momentum failure --- AdamW train loss collapses to zero while SGD stays pinned at the uniform-distribution value ln(99) as the weight norm decays 42 → 0.02.}]{fig4_negative_results.png}}
\caption{Left: sustained t\_grok increases monotonically with batch size
at p=97, wd=1.0 (1067 → 1314 → 1467 for bs 256 → 512 → 1024): minibatch
noise accelerates grokking. Right: the SGD-momentum failure --- AdamW
train loss collapses to zero while SGD stays pinned at the
uniform-distribution value ln(99) as the weight norm decays 42 → 0.02.}
\end{figure}

\subsection{5. Theory: status and
method}\label{theory-status-and-method}

\subsubsection{5.1 A descriptive
framework}\label{a-descriptive-framework}

Our working model (developed in \texttt{theory/}) treats grokking as
circuit selection under a complexity budget: t\_grok ∝ exp(L·μ/D) with
wd capping circuit complexity L, μ counting data-compatible circuits,
and D a noise/temperature term. A two-parameter reduction fits the
six-point α(K) curve with residuals 0.03--0.13. \textbf{We claim
descriptive adequacy only} (Section 5.2 explains why the framework
currently has zero confirmatory standing) --- except that H4's sign
(Section 4.4) already forced one revision: the noise term must act as
temperature (T ∝ lr/bs, higher is faster), not as signal-to-noise.

\subsubsection{5.2 An integrity
disclosure}\label{an-integrity-disclosure}

An earlier version of this framework published ``pre-registered''
predictions for the intermediate moduli, labeled as written before those
data were observed, and they passed their falsification test (all three
within the stated 0.10 line). During preparation of this paper we
reconstructed the fit --- no fit script had been saved --- and found:
(i) the document's stated parameters reproduce neither its predictions
nor its own anchors; (ii) an honest re-derivation from the anchors it
claims to have used \emph{fails} the falsification line at p=67; (iii)
the published predictions instead match, to ≤0.006 on all three points,
a fit using slope values that only exist with the later data; and (iv)
file timestamps show those data existed \textasciitilde14 hours before
the document. We treat the episode as HARKing, vacate the confirmatory
claim, and freeze the original document with the audit in its changelog
(\texttt{prereg/theory\_extension\_v0.1.md}; reconstruction:
\texttt{code/fit\_alpha\_K.py}; figure
\texttt{alpha\_vs\_K\_pred\_vs\_obs.png} shows the published
``predictions'' sitting on the contaminated fit). We publish the failure
because the alternative --- quietly dropping the claim --- is how fields
accumulate unreproducible theory.

\subsubsection{5.3 Replacement predictions, registered properly --- one
passed, one
failed}\label{replacement-predictions-registered-properly-one-passed-one-failed}

Two consequences of the temperature reading were pre-registered in
\texttt{prereg/PRE\_REGISTRATION\_P1\_P4.md}, timestamped by public
commit \texttt{320f7799} (2026-07-07), \textbf{before} their runs began
(33/33 runs clean, \texttt{results/phase9/}):

\begin{itemize}
\tightlist
\item
  \textbf{P1 --- CONFIRMED.} At fixed batch size, t\_grok decreases
  monotonically with learning rate (2567 → 1333 → 933 for lr
  5e-4/1e-3/2e-3), the direction the temperature reading T = lr/bs
  predicts. The magnitude at low lr exceeds the exploratory band (ratio
  1.93 vs predicted ≤1.55) --- direction right, low-lr slowdown stronger
  than modeled.
\item
  \textbf{P4 --- REFUTED, and this is the more informative outcome.} The
  wd exponent is \emph{not} batch-size-invariant: α = -0.636, -0.731,
  -0.802 at bs = 256, 512, 1024. The pre-registered separation
  \textbar α(256) − α(1024)\textbar{} ≤ 0.15 fails at 0.166. Worse for
  the model, better for the science: \textbf{\textbar α\textbar{}
  increases monotonically with batch size} --- larger batches make
  grokking time \emph{more} sensitive to weight decay. A separable
  temperature prefactor cannot produce this; weight decay and batch size
  interact.
\end{itemize}

Per the pre-registered decision table (P1 pass + P4 fail), the separable
temperature model is rejected and an interaction model is required
(developed as eq. (11) in \texttt{theory/}); its own sharpest test (does
α depend on lr as it does on bs?) is left for the next pre-registration.
We report the 1-of-2 hit rate exactly as the public commit obliges. The
value of P4 is that it was a real risk, taken publicly, that the theory
lost --- which is the only kind of test that can earn a theory credit
when it wins.

\subsection{6. Related work}\label{related-work}

Power et al.~(2022) established grokking on modular arithmetic and the
role of weight decay. Nanda et al.~(2023) identified the Fourier-circuit
mechanism at p=113 and proposed progress measures; our map is consistent
with circuit selection being the rate-limiting process. Liu et
al.~(2022) frame grokking as a phase transition; α(K) is naturally read
as a K-dependent critical exponent, though we make no formal claim.
Loshchilov \& Hutter (2019) motivate why coupled vs decoupled weight
decay behave differently, which Section 4.2 turns into a hard scope
boundary. The weight-norm collapse we observe is the SGD side of the
equilibrium analysis of Kosson et al.~(2023), who show that weight decay
balances update and decay only under specific optimizer dynamics;
D'Angelo et al.~(2023) likewise find that weight decay's role differs
qualitatively between SGD and Adam-class training. That transformers
train poorly under plain SGD is documented by Zhang et al.~(2019)
(heavy-tailed gradient noise) and Kunstner et al.~(2023) (sign-descent
behaviour, not noise, explains most of the gap); our contribution is the
concrete failure mode at these hyperparameters. The temperature reading
of Sections 4.4--5.3, T ∝ lr/bs, is the SGD noise scale of Mandt et
al.~(2017) and Smith \& Le (2017); the P4 interaction says this
prefactor is \emph{not} separable from weight decay in the grokking
regime. Adaptive-optimizer instabilities during grokking are studied by
Thilak et al.~(2022); theoretical accounts of the memorization-to-
generalization transition include lazy-to-rich dynamics (Kumar et al.,
2023), the early/late implicit-bias dichotomy (Lyu et al., 2023) and
circuit efficiency (Varma et al., 2023), whose competing-circuit picture
is the one our informal interpretation in Section 3 leans on.

\subsection{7. Limitations}\label{limitations}

\begin{enumerate}
\def\labelenumi{\arabic{enumi}.}
\tightlist
\item
  One task family (modular addition), one architecture, p ≤ 149. No
  universality claim.
\item
  N=3 for most cells; strict monotonicity of α(K) unresolved; the
  steep-end anchor (p=31, wd=0.01) is fragile (3/8 clean).
\item
  The theoretical framework is post-hoc (Section 5.2); its first genuine
  tests are P1/P4.
\item
  Depth-2 conclusions concern a different (non-grokking) regime; see
  4.3.
\item
  Compute limits: 200K-step budget makes some low-wd cells lower bounds.
\end{enumerate}

\subsection{8. Conclusion}\label{conclusion}

Weight decay's leverage over grokking is not a constant of the
phenomenon: it decays with task size, \textbar α\textbar{} falling from
0.93 to 0.58 across a 23× range of K. The effect survives initialization
and batch-size checks, is bounded to AdamW-class optimizers by a
mechanistically-explained SGD failure, and persists across a regime
boundary where grokking itself disappears. The accompanying theory is
offered as description only: of its two publicly pre-registered
predictions, the learning-rate direction held and batch-size invariance
of α failed, which rejects a separable temperature model and leaves a
weight-decay × batch-size interaction as the next thing to explain ---
and to pre-register.

\subsection{Reproducibility}\label{reproducibility}

All code, per-run CSVs (including outliers, DNFs, and the truncated run
from a mid-phase crash), auto-generated reports, and dated
pre-registration artifacts:
\texttt{github.com/DoctorMilesBennettDyson/grokking-alpha-k}. Hardware:
single RTX 3070 8 GB; PyTorch 2.4.1+cu118. Runners are idempotent;
re-running resumes.

\subsection{References}\label{references}

\begin{itemize}
\tightlist
\item
  Power, A., Burda, Y., Edwards, H., Babuschkin, I., Misra, V. (2022).
  \emph{Grokking: Generalization Beyond Overfitting on Small Algorithmic
  Datasets.} arXiv:2201.02177.
\item
  Nanda, N., Chan, L., Lieberum, T., Smith, J., Steinhardt, J. (2023).
  \emph{Progress Measures for Grokking via Mechanistic
  Interpretability.} ICLR 2023. arXiv:2301.05217.
\item
  Liu, Z., Kitouni, O., Nolte, N., Michaud, E., Tegmark, M., Williams,
  M. (2022). \emph{Towards Understanding Grokking: An Effective Theory
  of Representation Learning.} NeurIPS 2022. arXiv:2205.10343.
\item
  Loshchilov, I., Hutter, F. (2019). \emph{Decoupled Weight Decay
  Regularization.} ICLR 2019. arXiv:1711.05101.
\item
  Mandt, S., Hoffman, M. D., Blei, D. M. (2017). \emph{Stochastic
  Gradient Descent as Approximate Bayesian Inference.} JMLR 18.
  arXiv:1704.04289.
\item
  Smith, S. L., Le, Q. V. (2017). \emph{A Bayesian Perspective on
  Generalization and Stochastic Gradient Descent.} arXiv:1710.06451.
\item
  Zhang, J., Karimireddy, S. P., Veit, A., Kim, S., Reddi, S. J., Kumar,
  S., Sra, S. (2019). \emph{Why are Adaptive Methods Good for Attention
  Models?} arXiv:1912.03194.
\item
  Kunstner, F., Chen, J., Lavington, J. W., Schmidt, M. (2023).
  \emph{Noise Is Not the Main Factor Behind the Gap Between SGD and Adam
  on Transformers, but Sign Descent Might Be.} arXiv:2304.13960.
\item
  Kosson, A., Messmer, B., Jaggi, M. (2023). \emph{Rotational
  Equilibrium: How Weight Decay Balances Learning Across Neural
  Networks.} arXiv:2305.17212.
\item
  D'Angelo, F., Andriushchenko, M., Varre, A., Flammarion, N. (2023).
  \emph{Why Do We Need Weight Decay in Modern Deep Learning?}
  arXiv:2310.04415.
\item
  Thilak, V., Littwin, E., Zhai, S., Saremi, O., Paiss, R., Susskind, J.
  (2022). \emph{The Slingshot Mechanism: An Empirical Study of Adaptive
  Optimizers and the Grokking Phenomenon.} arXiv:2206.04817.
\item
  Kumar, T., Bordelon, B., Gershman, S. J., Pehlevan, C. (2023).
  \emph{Grokking as the Transition from Lazy to Rich Training Dynamics.}
  arXiv:2310.06110.
\item
  Lyu, K., Jin, J., Li, Z., Du, S. S., Lee, J. D., Hu, W. (2023).
  \emph{Dichotomy of Early and Late Phase Implicit Biases Can Provably
  Induce Grokking.} arXiv:2311.18817.
\item
  Varma, V., Shah, R., Kenton, Z., Kramár, J., Kumar, R. (2023).
  \emph{Explaining grokking through circuit efficiency.}
  arXiv:2309.02390.
\end{itemize}

\end{document}
